title: ``Отчёт по лабораторной работе 1'' subtitle: ``Подготовка
лабораторного стенда'' author: ``Талебу Тенке Франк Устон'' institute:
``Российский университет дружбы народов'' email: ``1032224534@pfur.ru''
date: ``2026-03-23'' ---

\subsection{Generic otions}\label{generic-otions}

lang: ru-RU toc-title: ``Содержание''

\subsection{Bibliography}\label{bibliography}

bibliography: bib/cite.bib csl: pandoc/csl/gost-r-7-0-5-2008-numeric.csl

\subsection{Pdf output format}\label{pdf-output-format}

toc: true \# Table of contents toc-depth: 2 lof: true \# List of figures
lot: true \# List of tables fontsize: 12pt linestretch: 1.5 papersize: a
documentclass: scrreprt \#\# I18n polyglossia polyglossia-lang: name:
russian options: - spelling=modern - babelshorthands=true
polyglossia-otherlangs: name: english \#\# I18n babel babel-lang:
russian babel-otherlangs: english \#\# Fonts mainfont: IBM Plex Serif
romanfont: IBM Plex Serif sansfont: IBM Plex Sans monofont: IBM Plex
Mono mathfont: STIX Two Math mainfontoptions:
Ligatures=Common,Ligatures=TeX,Scale=0.94 romanfontoptions:
Ligatures=Common,Ligatures=TeX,Scale=0.94 sansfontoptions:
Ligatures=Common,Ligatures=TeX,Scale=MatchLowercase,Scale=0.94
monofontoptions: Scale=MatchLowercase,Scale=0.94,FakeStretch=0.9
mathfontoptions: \#\# Biblatex biblatex: true biblio-style:
``gost-numeric'' biblatexoptions: - parentracker=true - backend=biber -
hyperref=auto - language=auto - autolang=other* - citestyle=gost-numeric
\#\# Pandoc-crossref LaTeX customization figureTitle: ``Рис.''
tableTitle: ``Таблица'' listingTitle: ``Листинг'' lofTitle: ``Список
иллюстраций'' lotTitle: ``Список таблиц'' lolTitle: ``Листинги'' \#\#
Misc options indent: true header-includes: -

\usepackage{indentfirst}

\begin{itemize}
\item
  \usepackage{float}

  \section{keep figures where there are in the
  text}\label{keep-figures-where-there-are-in-the-text}

  \begin{itemize}
  \item ~
    \subsection{\texorpdfstring{\floatplacement{figure}{H} \# keep
    figures where there are in the
    text}{ \# keep figures where there are in the text}}\label{keep-figures-where-there-are-in-the-text-1}
  \end{itemize}
\end{itemize}

\section{Цель
работы}\label{ux446ux435ux43bux44c-ux440ux430ux431ux43eux442ux44b}

Целью данной работы является приобретение практических навыков установки
Rocky Linux на виртуальную машину с помощью инструмента Vagrant.

\section{Выполнение}\label{ux432ux44bux43fux43eux43bux43dux435ux43dux438ux435}

\subsection{Развёртывание лабораторного стенда на ОС
Windows}\label{ux440ux430ux437ux432ux451ux440ux442ux44bux432ux430ux43dux438ux435-ux43bux430ux431ux43eux440ux430ux442ux43eux440ux43dux43eux433ux43e-ux441ux442ux435ux43dux434ux430-ux43dux430-ux43eux441-windows}

\begin{enumerate}
\def\labelenumi{\arabic{enumi}.}
\item
  Подготовлена рабочая среда на локальном компьютере под управлением
  Windows.\\
  Установлены и настроены VirtualBox и Vagrant. После установки система
  была перезагружена.
\item
  В рабочем каталоге проекта созданы необходимые файлы для
  автоматического развертывания --- 01-user.sh, 01-hostname.sh и
  Vagrantfile.

  \begin{figure}
  \centering
  
\includegraphics[width=0.8\linewidth,height=\textheight,keepaspectratio]{1.png}
  \caption{Скрипт 01-user.sh}\label{fig:001}
  \end{figure}

  \begin{figure}
  \centering
  \includegraphics[width=0.8\linewidth,height=\textheight,keepaspectratio]{2.png}
  \caption{Скрипт 01-hostname.sh}\label{fig:002}
  \end{figure}

  \begin{figure}
  \centering
  \includegraphics[width=0.8\linewidth,height=\textheight,keepaspectratio]{3.png}
  \caption{Фрагмент Vagrantfile с описанием provision}\label{fig:003}
  \end{figure}
\item
  Выполнена регистрация box-файла в Vagrant с помощью команды:

  vagrant box add rockylinux10
  vagrant-virtualbox-rockylinux10-x86\_64.box

  \begin{figure}
  \centering
  \includegraphics[width=0.8\linewidth,height=\textheight,keepaspectratio]{4.png}
  \caption{Добавление box-файла в Vagrant}\label{fig:004}
  \end{figure}
\item
  Выполнен запуск виртуальной машины Server:

  vagrant up server

  \begin{figure}
  \centering
  \includegraphics[width=0.8\linewidth,height=\textheight,keepaspectratio]{5.png}
  \caption{Запуск виртуальной машины server}\label{fig:005}
  \end{figure}
\item
  После успешного запуска на экране входа в систему появились
  пользователи vagrant и созданный скриптом akabrikosov.

  \begin{figure}
  \centering
  \includegraphics[width=0.8\linewidth,height=\textheight,keepaspectratio]{6.png}
  \caption{Окно входа в систему Rocky Linux}\label{fig:006}
  \end{figure}
\item
  Проверено, что при входе в систему отображается корректное имя хоста
  server.akabrikosov.net.

  \begin{figure}
  \centering
  \includegraphics[width=0.8\linewidth,height=\textheight,keepaspectratio]{7.png}
  \caption{Проверка имени хоста в терминале}\label{fig:007}
  \end{figure}
\item
  Подключение к серверу через консоль выполнено командой:

  vagrant ssh server

  После ввода пароля осуществлён вход и произведена проверка возможности
  переключения на пользователя akabrikosov:

  su - akabrikosov

  Затем выполнен выход из системы.

  \begin{figure}
  \centering
  \includegraphics[width=0.8\linewidth,height=\textheight,keepaspectratio]{8.png}
  \caption{Подключение к серверу через SSH}\label{fig:008}
  \end{figure}
\end{enumerate}

\section{Вывод}\label{ux432ux44bux432ux43eux434}

В результате работы был развернут лабораторный стенд на базе Vagrant и
VirtualBox в среде Windows.\\
Созданы и протестированы сценарии автоматического развертывания
виртуальной машины с использованием скриптов 01-user.sh и
01-hostname.sh.\\
Проверена корректность работы созданного пользователя, hostname и
подключения через SSH.

\section{Контрольные
вопросы}\label{ux43aux43eux43dux442ux440ux43eux43bux44cux43dux44bux435-ux432ux43eux43fux440ux43eux441ux44b}

\begin{enumerate}
\def\labelenumi{\arabic{enumi}.}
\item
  Для чего предназначен Vagrant\\
  Vagrant предназначен для автоматизации развертывания и управления
  виртуальными машинами.\\
  Он упрощает настройку окружений для разработки и тестирования,
  обеспечивая единообразие конфигураций.
\item
  Что такое box-файл? В чём назначение Vagrantfile\\
  Box-файл --- это шаблон готовой виртуальной машины, содержащий
  операционную систему и базовые настройки.\\
  Vagrantfile --- это конфигурационный файл, описывающий параметры
  виртуальных машин: используемый box, настройки сети, ресурсы, а также
  сценарии инициализации (provision).
\item
  Приведите описание и примеры вызова основных команд Vagrant

  \begin{itemize}
  \tightlist
  \item
    \texttt{vagrant\ init} --- создаёт новый Vagrantfile в текущем
    каталоге.\\
  \item
    \texttt{vagrant\ box\ add\ \textless{}имя\textgreater{}\ \textless{}путь\ к\ box\textgreater{}}
    --- добавляет образ виртуальной машины в локальное хранилище.\\
  \item
    \texttt{vagrant\ up} --- запускает и создаёт виртуальную машину.\\
  \item
    \texttt{vagrant\ ssh} --- подключается к запущенной машине по SSH.\\
  \item
    \texttt{vagrant\ halt} --- корректно останавливает виртуальную
    машину.\\
  \item
    \texttt{vagrant\ destroy} --- полностью удаляет виртуальную
    машину.\\
  \item
    \texttt{vagrant\ status} --- показывает текущее состояние всех машин
    проекта.
  \end{itemize}
\item
  Дайте построчные пояснения содержания файлов vagrant-rocky.pkr.hcl,
  ks.cfg, Vagrantfile, Makefile

  \begin{itemize}
  \tightlist
  \item
    \textbf{vagrant-rocky.pkr.hcl} --- конфигурация для утилиты Packer,
    определяющая процесс сборки box-файла (указывается источник ISO,
    параметры установки и скрипты автоматизации).\\
  \item
    \textbf{ks.cfg} --- файл автоматической установки системы
    (kickstart), в котором задаются параметры установки ОС: язык,
    разделы диска, учетные данные, пакеты.\\
  \item
    \textbf{Vagrantfile} --- основной сценарий Vagrant, описывающий
    создание виртуальных машин (server, client), их параметры сети,
    провижнинг и запуск вспомогательных скриптов.\\
  \item
    \textbf{Makefile} --- содержит набор команд для автоматизации часто
    выполняемых действий, таких как сборка, запуск или уничтожение
    виртуальных машин.
  \end{itemize}
\end{enumerate}
